\BourbakiExerciseGroup

\begin{enumerate}
\def\labelenumi{\arabic{enumi})}
\item
  Show that every bounded half-open interval of \(\mathbf{R}\) is homeomorphic to every unbounded closed interval of \(\mathbf{R}\) . If \(a < b\) , then no two of the three intervals \(]a, b[, [a, b[, [a, b]]\) are homeomorphic.
\item
  Show that the mapping \(x \to x / (1 + |x|)\) is a homeomorphism of \(\mathbf{R}\) onto the open interval \(] - 1, 1[\) .
\item
  Let I denote the interval \([0, 1]\) of \(\mathbf{R}\) , and let \(f\) be a homeomorphism of I onto itself such that \(f(0) = 0\) and \(f(1) = 1\) . Show that there is a continuous mapping \(g\) of \(\mathbf{I} \times \mathbf{I}\) into I such that: (i) for each \(x \in \mathbf{I}\) , \(g(x, 0) = x\) and \(g(x, 1) = f(x)\) ; (ii) for each \(y \in \mathbf{I}\) the partial mapping \(x \to g(x, y)\) is a homeomorphism of I onto itself such that \(g(0, y) = 0\) and \(g(1, y) = 1\) .
\item
  Show that the uniformity induced on \(\mathbf{R}\) by the unique uniformity of \(\overline{\mathbf{R}}\) is strictly coarser than the additive uniformity of \(\mathbf{R}\) .
\item
  As the point \((x, y)\) tends to \((+\infty, -\infty)\) while remaining in \(\mathbf{R} \times \mathbf{R}\) , shows that the set of cluster points of the function \(x + y\) is \(\overline{\mathbf{R}}\) .
\end{enumerate}

Likewise, as \((x,y)\) tends to \((0,+\infty)\) while remaining in \(R\times R\) , the set of cluster points of xy is \(\overline{R}\) .

\begin{enumerate}
\def\labelenumi{\arabic{enumi})}
\setcounter{enumi}{5}
\item
  Show that every rational function \(\mathbf{P}(x)/\mathbf{Q}(x)\) with coefficients in R, defined at all points x of R such that \(\mathbf{Q}(x)\neq0\) , can be extended by continuity to the points \(+\infty\) and \(-\infty\) (taking its values in \(\overline{R}\) ).
\item
  Let X be a linearly ordered set.
\end{enumerate}

\begin{enumerate}
\def\labelenumi{\alph{enumi})}
\item
  Let \(X_{1}\) be the completion of X (S et Theory, Chapter III, § 1, Exercise 15). \(X_{1}\) is a linearly ordered set whose elements are the subsets A of X such that; (i) if \(x \in A\) and \(y \leqslant x\) , then \(y \in A\) ; (ii) if A has a least upper bound in X, this bound belongs to A {[}which is equivalent to saying that A is closed in \(\mathcal{G}_{0}(X)\) {]}; the elements of \(X_{1}\) are ordered by inclusion. \(X_{1}\) is compact with respect to the topology \(\mathcal{G}_{0}(X_{1})\) {[}§ 2, Exercise 6 a){]}. The mapping \(x \to ]\leftarrow,x]\) of X into \(X_{1}\) is strictly order-preserving, and if we identify X with its image under this mapping, then X is dense in \(X_{1}\) and \(\mathcal{G}_{0}(X)\) is induced by \(\mathcal{G}_{0}(X_{1})\) . \(X_{1}\) is connected if and only if X is without gaps (§ 2, Exercise 7).
\item
  For each cardinal c, give an example of a linearly ordered set X which is connected in the topology \(\mathfrak{G}_{0}(\mathbf{X})\) and is such that every fundamental system of neighbourhoods of any point of X has cardinal \(\geqslant c\) {[}use a) and the method of §1, Exercise 4 b){]}.
\item
  With the notation of \(a\) , let \(\mathbf{X}_{2}\) be the subset of the lexicographic product \(\mathbf{X}_{1} \times \left\{-\mathbf{i}, 0, \mathbf{i}\right\}\) which is the complement of the set consisting of: (i) the points \((x, 0)\) where \(x \notin \mathbf{X}\) ; (ii) the points \((x, 1)\) where \(x \in \mathbf{X}\) and the set of all \(y > x\) in \(\mathbf{X}\) has a smallest element; (iii) the points \((x, -1)\) where \(x \in \mathbf{X}\) and the set of all \(y < x\) in \(\mathbf{X}\) has a greatest element. Show that \(\mathbf{X}_{2}\) , endowed with the topology \(\mathfrak{C}_{0}(\mathbf{X}_{2})\) , is compact and totally disconnected, that \(\mathbf{X}\) can be identified, by means of the strictly order-preserving mapping \(x \to (x, 0)\) , with a dense subset of \(\mathbf{X}_{2}\) , and that the topology induced by \(\mathfrak{C}_{0}(\mathbf{X}_{2})\) on \(\mathbf{X}\) is discrete.
\item
  Let \(X'\) be a linearly ordered set containing X which induces the given order structure on X, is compact in the topology \(\mathcal{C}_{0}(\mathbf{X}')\) and is such that X is dense in \(X'\) with respect to this topology. Show that there is a continuous order-preserving surjective mapping \(f: X' \to X_{1}\) and a continuous order-preserving surjective mapping \(g: X_{2} \to X'\) , both of which induce the identity mapping on X.
\end{enumerate}

